\section{Background and Motivation}\label{sec:background}

\subsection{Python Bytecode and Code Objects}

Python source follows a compilation pipeline from source text to abstract syntax tree, code object, and bytecode execution.  A code object contains the bytecode instruction stream together with constants, names, local-variable metadata, free and cell variables, flags, line and position information, and, in newer CPython versions, exception-table metadata.  CPython serializes code objects using the \texttt{marshal} format and stores import caches as \texttt{.pyc} files under \texttt{\_\_pycache\_\_}.  These artifacts are version-sensitive: opcode sets, instruction layouts, stack effects, and metadata formats change across Python releases.

Although \texttt{.pyc} files are often described as caches, they are executable representations consumed by the interpreter.  Python can import bytecode-only modules, load code objects from custom loaders, execute unmarshalled code through \texttt{exec}, and ship compiled payloads inside application bundles.  Therefore, bytecode is not merely a performance optimization; it is part of the attack and analysis surface.  This matters even when a package also contains source, because the source file and the bytecode artifact may differ in version, path, timestamp, or content.

\subsection{Security Relevance of Bytecode Artifacts}

Bytecode artifacts arise in several security-relevant settings.  Package distributions may accidentally include cached bytecode.  Deployment systems may ship bytecode without source.  Bundlers and release systems may embed compiled modules.  Malware can hide payloads in marshalled code objects or source-less modules.  Security scanners may encounter all of these artifacts while analyzing wheels, source distributions, containers, virtual environments, or release archives.

The challenge is that source-level visibility is no longer guaranteed.  A nearby \texttt{.py} file may be absent, stale, or unrelated to the \texttt{.pyc} that will be loaded.  Decompilers may recover readable source for common cases, but they are not a complete substitute for bytecode analysis.  A failed decompilation does not imply benign behavior, and a successful decompilation does not guarantee semantic equivalence after recompilation.  For this reason, source presence, source matching, tool recoverability, and runtime behavior must be measured as separate properties.

\subsection{Motivating Example}\label{sec:motiexample}

Consider a package that contains a benign-looking module \texttt{loader.py} and a bytecode artifact embedded as bytes.  The source-level loader may only call \texttt{marshal.loads} and \texttt{exec}; the behavior of interest is inside the recovered code object.  A scanner that indexes only \texttt{.py} syntax sees a generic dynamic execution site, while the interpreter eventually consumes concrete bytecode.  If the embedded code object targets a different CPython version, uses an opcode unsupported by a decompiler, or has control-flow metadata that source recovery cannot reconstruct, source-level analysis becomes incomplete.

This example illustrates the two dimensions studied in this paper.  First, bytecode artifacts can create a visibility gap between what package security tools inspect and what the runtime can execute.  Second, bytecode artifacts can exercise interpreter paths that are difficult to reach from ordinary source, making runtime robustness results hard to interpret without a tool-bounded source-reproduction analysis.
